\section{Feedback Learning}
\label{sec:feedback_learning}

Feedback learning studies how completed trajectories change behavior on future tasks. It differs from adaptive retrieval within a single task: the defining update persists beyond the current episode. The retained object may be a reflection, an experience record, a memory item, a reusable procedure, a skill, retriever supervision, or an updated search policy.

\subsection{Feedback Collection}
\label{subsec:feedback_collection}
% Cover outcome rewards, process rewards, verifier signals, human feedback,
% tool feedback, environment feedback, and automatic critics.

\subsection{Trajectory Diagnosis}
\label{subsec:trajectory_diagnosis}
% Cover failure attribution, reflection, comparison of successful and failed
% trajectories, credit assignment, and diagnosis granularity.

\subsection{Experience Memory}
\label{subsec:experience_memory}
% Cover episodic experience stores, distilled lessons, retrieval of prior
% trajectories, memory consolidation, and interference or staleness.

\subsection{Skill Extraction}
\label{subsec:skill_extraction}
% Cover executable skills, textual procedures, skill libraries, automatic
% skill revision, skill routing, and compositional reuse.

\subsection{Retrieval and Policy Improvement}
\label{subsec:policy_improvement}
% Cover reinforcement-learned search policies, process-supervised retrievers,
% retriever updates from downstream trajectories, and continual adaptation.

\paragraph{Inclusion rule.}
General memory or skill systems are included only when their retained experience directly improves retrieval-grounded multi-hop behavior. A long-term memory module that merely stores user facts, or an embodied skill unrelated to information acquisition, is outside the primary scope.
